\documentclass[11pt]{article}
\usepackage[margin=1in]{geometry}
\usepackage{amsmath,amssymb,amsthm}
\usepackage{hyperref}
\usepackage{booktabs}
\usepackage{graphicx}
\hypersetup{hidelinks}

\newtheorem{theorem}{Theorem}
\newtheorem{lemma}[theorem]{Lemma}
\newtheorem{corollary}[theorem]{Corollary}
\newtheorem{proposition}[theorem]{Proposition}
\theoremstyle{definition}
\newtheorem{definition}{Definition}

\setlength{\parskip}{4pt plus 1pt minus 1pt}
\setlength{\parindent}{0pt}

\title{The Ratchet:\\Evolution and Cancer as Opposing Failures\\of the Same Mechanism}
\author{Patrick D.\ McCarthy}
\date{}

\begin{document}
\maketitle

\begin{abstract}
This paper completes the series. Papers~1--4 derive the structural
constraints of bounded context. Papers~5--8 validate them in cancer
genomics. Paper~9 shows their instantiation in a production system.
Paper~10 shows their instantiation in the genome. This paper names
the mechanism that connects them.

A knowledge graph under bounded context can fail in two directions.
Forward: novel signal exceeds the graph's capacity to encode it, and
the system must restructure. This is evolution. Backward: governance
degrades and the system loses the capacity to maintain its own
structure. This is cancer. The mechanism is the same. The direction
differs. We call it the ratchet.

The ratchet has a characteristic shape: the sigmoid. Two scaling
classes ($O(n^2)$ and $O(\log n)$) compete. The inferior class
outperforms at small scale. At a crossover point, it fails
catastrophically. Every system transition in the series follows
this curve. We derive its shape from the bounded context constraint
and show that cancer, evolution, aging, and the current crisis in
AI scaling are all instances of the same crossover.

The genome and the software industry arrived at identical
architectures not by analogy but by necessity. Bounded context
admits only one family of solutions. 3.8 billion years of selection
and 60 years of engineering found the same one.
\end{abstract}

% ------------------------------------------------------------------
\section{Introduction}
\label{sec:intro}
% ------------------------------------------------------------------

The series began with a constraint: active context is finite. Every
paper since has been a consequence of that constraint. This paper is
the last consequence.

Paper~1 proved that bounded context requires graph structure.
Paper~2 proved that graphs under survival pressure escalate.
Paper~3 proved that knowledge and action cannot be separated.
Paper~4 proved that encoding permanence must track confidence.
Papers~5--8 showed that cancer data obeys these constraints.
Paper~9 showed a production system discovering them.
Paper~10 showed the genome implementing them.

This paper asks: what happens when the constraints are violated?

The answer is the ratchet. It clicks forward (evolution) or backward
(cancer). Both are failures of the same mechanism. The mechanism is
governance under bounded context. The direction depends on which
force dominates: novelty from outside the graph or degradation from
within it.

% ------------------------------------------------------------------
\section{The Graph Changes First}
\label{sec:graph-first}
% ------------------------------------------------------------------

The standard model of molecular biology places DNA at the center:
DNA $\to$ RNA $\to$ protein. Changes begin in the sequence and
propagate outward. This paper inverts the arrow.

The cell is a running knowledge graph. Its Meta/Action/Knowledge
state (Paper~1) is maintained in protein concentrations, metabolite
levels, signaling cascades, and chromatin accessibility. DNA is the
serialized projection of this graph (Paper~10). Epigenetic marks are
the write-ahead log. The sequence is the commit log.

Changes begin in the running graph, not in the sequence. A cell
responds to a signal. Its M/A/K state shifts. If the shift is small,
the epigenetic layer (the write-ahead log) absorbs it. The scratch
pad handles the transaction. No sequence change is needed.

If the shift exceeds the scratch pad's capacity, the change must be
committed to the sequence. This is a one-way door. The scratch pad
(epigenetics) is lossy: transgenerational epigenetic inheritance fades
after 3--5 generations \cite{Heard2014} because the WAL has
approximately one layer of error correction (Paper~10). The commit log (DNA) persists
indefinitely because it has seven.

The one-way door is the ratchet. Small changes are absorbed and may
wash out. Large changes, if they survive long enough to exceed the
scratch pad, commit. Once committed, they do not uncommit. The
sequence ratchets forward.

\subsection{Evidence: epigenetics precedes genetics}

This inversion is not novel. The epigenetic progenitor model
\cite{Feinberg2006} and the Waddington landscape tradition
established that epigenetic changes precede genetic ones in
carcinogenesis. The contribution here is not the observation but
the formalization: the WAL/commit-log architecture from Paper~10
explains \emph{why} this ordering is necessary, not merely that
it is observed.

If the graph changes first, epigenetic changes should precede
sequence changes. They do. DNA methylation changes
are detectable in pre-malignant tissue years before driver mutations
appear \cite{Feinberg2006}. The scratch pad fills before the commit
log corrupts. CHIP (Paper~10) is this sequence in slow motion:
epigenetic maintenance fails first, and sequence-level driver
mutations follow years to decades later \cite{Jaiswal2014}.

Field cancerization---the observation that histologically normal
tissue adjacent to tumors carries epigenetic but not genetic
alterations \cite{Slaughter1953}---is the spatial version of the
same sequence. The graph is already changed. The commit has not yet
happened.

\subsection{One mutation can spawn an organ}

A single germline mutation does not build an organ. It changes the
decompression instructions (Paper~10, selective decompression).
The complexity was already latent in the projection. The mutation
unlocks a subgraph that was never previously decompressed.

This resolves the paradox of how a point mutation can produce a
complex phenotype. The phenotype is not encoded in the mutation.
It is encoded in the graph. The mutation changes which part of the
graph gets traversed. A single switch in a routing table redirects
traffic to an entire subnetwork that was always present but never
addressed.

% ------------------------------------------------------------------
\section{The Ratchet}
\label{sec:ratchet}
% ------------------------------------------------------------------

The ratchet clicks in two directions.

\subsection{Forward: evolution}

A novel signal arrives from outside the graph. Environmental change.
A new pathogen. A new food source. The signal cannot be handled by
the current M/A/K state. The system must restructure.

If the restructuring is small, epigenetics absorbs it. The population
adapts without sequence change. This is phenotypic plasticity.

If the restructuring exceeds the scratch pad, a sequence change is
needed. The mutation that enables the restructuring is selected for.
It commits. The ratchet clicks forward.

Forward ratcheting builds larger, more complex graphs. More nodes.
More edges. More layers of governance (Paper~10). The organism gains
capacity. It can handle more signal. It can maintain more context.
It moves rightward on the scaling curve.

But forward ratcheting has a cost: every new node in the graph is a
new node that must be governed. The governance burden scales as
$\log(n)$ (Paper~10), but it scales. The organism becomes more
capable and more vulnerable simultaneously. More context capacity
means more context to maintain. More governance layers means more
layers that can fail.

\subsection{Backward: cancer}

Governance degrades. A repair layer loses confidence. A checkpoint
gene mutates. The WAL replay machinery fails (CHIP). The system
cannot maintain its own structure.

The cell does not ``decide'' to become cancerous. It loses the
capacity to maintain its position in the graph. Its M/A/K state
drifts. It stops responding to coordination signals. It stops
respecting resource boundaries. It reverts to the strategy that
preceded governance: replicate, consume, compete.

This is not reversion to an ancestral form. The bacterial mutation
rate ($\sim 10^{-6}$) is not something cancer cells accelerate to.
It is the base rate. It was always there. The seven governance layers
suppressed it. As layers are shed, the base rate re-emerges
(Paper~10). The cell is not going backward. The governance that was
holding it forward is falling away.

Cancer is the ratchet clicking backward: not because the sequence
reverts, but because the governance that made the sequence
meaningful degrades. The cell defects from the graph to the swarm.
From Strategy~2 (maximize context) to Strategy~1 (minimize context
need). From $O(\log n)$ back to $O(n^2)$.

\subsection{The asymmetry}

Forward ratcheting requires something rare: a mutation that enables
a novel, beneficial restructuring of the graph. It requires signal
from outside the system.

Backward ratcheting requires something common: degradation of any
one of $k$ governance layers. It requires only entropy.

Evolution is hard. Cancer is easy. The forward ratchet clicks rarely.
The backward ratchet clicks constantly. The organism exists in the
interval between them: building complexity that entropy is always
trying to dismantle. This is the vise from Paper~1. Novelty from
above, degradation from below. The graph is the structure that
survives between them.

% ------------------------------------------------------------------
\section{The Sigmoid}
\label{sec:sigmoid}
% ------------------------------------------------------------------

\subsection{Two scaling classes}

Paper~1 identified two strategies for bounded context:
\begin{enumerate}
\item \textbf{Minimize context need.} Stay flat. Stay fungible.
   Horizontal scaling. $O(n^2)$ search over a small domain.
   Bacteria, BitTorrent, transformers.
\item \textbf{Maximize context capacity.} Build hierarchy. Build
   a graph. Vertical scaling. $O(\log n)$ indexed retrieval.
   Multicellular life, knowledge graphs, compiled code.
\end{enumerate}

At small scale, Strategy~1 outperforms. When $|K| \ll C_n$,
exhaustive search is faster than indexed retrieval because the
constant overhead of index maintenance dominates the exponent.
The dangerous property of the inferior scaling class is not that it
fails but that it succeeds, conspicuously, until a crossover it
cannot see.

\subsection{The crossover}

The crossover occurs where the total cost of Strategy~1 exceeds
the total cost of Strategy~2 (governance overhead included).
Below this point, Strategy~1 dominates. Above it, Strategy~2
dominates.

For any two functions $f(n) = an^2$ and $g(n) = b \log n + c$,
the crossover point $n^*$ satisfies $a(n^*)^2 = b \log n^* + c$.
The ratio $f/g$ diverges monotonically. It is not itself a sigmoid.
But it does not need to be.

Every real system has a resource bound $R$. The context window is
finite. The blood supply is finite. The budget is finite. The cost
curve $f(n) = an^2$ exists on a domain bounded by $R$. Below $R$,
the system absorbs the quadratic cost. At $R$, the system fails.
The transition between ``absorbs'' and ``fails'' is bounded between
0 and 1. It has the shape of a sigmoid: flat for a long time, steep
at the transition, flat again after failure.

The dangerous property is not the shape. It is the scale. The
sigmoid is so large---the upper bound so far from the current
operating point---that from below, the curve looks logarithmic.
The system sees steady improvement. It extrapolates forward and
sees more of the same. The inflection point is real but it is
beyond the observable horizon. By the time the curvature becomes
visible, the system is past it. This is why the transition appears
sudden. It is not sudden. It was accumulating on the part of the
curve that looked flat because the sigmoid was too big to see.

\subsection{Every system transition follows this curve}

\begin{center}
\small
\begin{tabular}{lllll}
\toprule
\textbf{Transition} & \textbf{Strategy 1} & \textbf{Strategy 2} & \textbf{Crossover signal} \\
\midrule
Unicellular $\to$ multicellular & Replicate, compete & Coordinate, specialize & Colony size \\
Prokaryote $\to$ eukaryote & Flat genome & Compartmented nucleus & Genome size \\
Innate $\to$ adaptive immunity & Fixed receptors & V(D)J recombination & Pathogen diversity \\
Scripted $\to$ learned behavior & Reflexes & Neural plasticity & Environment variance \\
Agent $\to$ script (Paper~9) & Full context per call & Cached determinism & Uncertainty $< \tau$ \\
Context stuffing $\to$ graph & $O(n^2)$ token cost & $O(\log n)$ traversal & Context window size \\
\bottomrule
\end{tabular}
\end{center}

In every case, the system begins with Strategy~1 because it is
simpler and cheaper at small scale. When the domain exceeds the
crossover, the system either transitions to Strategy~2 or dies.
There is no third option. The scaling math does not negotiate.

\subsection{Cancer as reverse crossover}

Cancer reverses the transition. A cell that has been operating under
Strategy~2 (coordinated, governed, part of the tissue graph) loses
enough governance to fall below the crossover. It reverts to
Strategy~1: replicate, compete, consume.

This reversion succeeds locally. The cancer cell outcompetes its
neighbors because it is not paying the governance overhead. It is
faster, less constrained, more resource-hungry. For the same reason
that Strategy~1 outperforms at small scale. The cancer cell is
operating below the crossover where $O(n^2)$ wins.

But the organism is above the crossover. The cancer cell's local
success is the organism's systemic failure. The noisy neighbor
consumes bounded resources (Paper~10). The system that vertically
scaled cannot tolerate a component that horizontally scales
without authorization.

% ------------------------------------------------------------------
\section{The Mutation Rate Signature}
\label{sec:mutation-rate}
% ------------------------------------------------------------------

The mutation rate of a cell is a direct readout of its remaining
governance (Paper~10, Table~1 \cite{McCarthy2026j}). Normal human
cells: $\sim 5 \times 10^{-10}$ (all seven layers). MSI-high
cancers: $\sim 10^{-7}$ (MMR lost). Hypermutated cancers:
$\sim 10^{-5}$ (multiple layers lost). The progression is not
``speeding up.'' It is shedding layers. The base rate was always
there. Cancer cells have lost the deceleration.

The convergence between late-stage cancer rates and bacterial rates
($\sim 10^{-6}$ to $10^{-3}$) is the prediction. Both represent the
base rate with minimal governance. The difference: bacteria evolved
to operate there. Cancer cells are operating there inside a system
that requires $10^{-10}$. The rate is the same. The context is
catastrophically different.

% ------------------------------------------------------------------
\section{Aging as Ratchet Drag}
\label{sec:aging}
% ------------------------------------------------------------------

If the forward ratchet is evolution and the backward ratchet is
cancer, aging is the friction in between: the slow accumulation of
governance degradation that has not yet crossed the threshold into
frank disease.

Paper~10 established the asymmetry: DNA (the data) has seven layers
of error correction. Epigenetic marks (the index) have approximately
one. Every cell division copies the data with high fidelity and the
index with low fidelity. Over a lifetime of divisions, the data
remains largely intact. The index degrades.

Aging is index degradation. The cell's data is fine. Its ability to
\emph{find} the right data---to maintain the right genes in the right
activation state in the right tissue---declines. Epigenetic drift is
the index losing coherence. Cellular senescence is the system
blocking writes when the WAL buffer is full (Paper~10). Stem cell
exhaustion is the system running out of instances that still have a
coherent index.

CHIP is the canary. DNMT3A, TET2, ASXL1---the index maintenance
machinery---fail first because they have the least redundancy
\cite{Jaiswal2014}. The data repair systems (proofreading, MMR, BER)
are still running. The checkpoints (p53) are still running. But the
WAL replay is failing, and every subsequent cell division loses more
index state.

The ratchet does not click backward all at once. It drags. Each
cell division with slightly worse index fidelity is a fraction of
a click. Over decades, the accumulated drag reaches a threshold
where some cell's governance falls below the crossover. That cell
reverts to Strategy~1. That is cancer. The ratchet has clicked.

% ------------------------------------------------------------------
\section{The Forward Ratchet in Software}
\label{sec:software}
% ------------------------------------------------------------------

The same mechanism operates in the systems described in Paper~9.

A software system begins simple: flat architecture, exhaustive
search, full context per query. Strategy~1. At small scale, this
outperforms. The system is fast, simple, and easy to reason about.

As the domain grows, the system hits the crossover. Context windows
fill. Response times degrade quadratically. The system must
restructure: build a knowledge graph, factor shared structure,
transition predictable actions to scripts. Strategy~2. The ratchet
clicks forward.

But forward ratcheting creates governance overhead. The knowledge
graph must be maintained. The scripts must be validated. The
transition thresholds must be monitored. Every abstraction is a
new node that must be governed. Paper~9's forever tax is the
governance cost of the forward ratchet.

And the backward ratchet operates too. Technical debt is governance
degradation. An abstraction that was correct becomes stale as the
domain shifts. A script that was validated against one input class
encounters a new class. A knowledge graph edge that was true becomes
false. Each degradation reduces the system's governance capacity.
Enough degradation and the system reverts: the team abandons the
graph, goes back to context stuffing, pays the $O(n^2)$ cost because
the governance overhead has become more expensive than the
alternative.

The software industry cycles between these states: build structure,
let it degrade, tear it down, build new structure. The ratchet clicks
forward when someone invests in governance. It drags backward when
no one does. The cycle is not a failure of engineering. It is the
ratchet operating on a shorter timescale than biological evolution
but under the same constraints.

Large language models are undergoing the transition now. The context
window is the bounded context. The quadratic attention cost is
$O(n^2)$. Retrieval-augmented generation, tool use, agent
architectures---these are the $O(\log n)$ infrastructure being built
under scaling pressure.

The parallel to the K-Pg extinction is precise. Dinosaurs at the
boundary occupied three positions on the sigmoid:

\begin{enumerate}
\item \textbf{Birds.} Had already begun building $O(\log n)$
   infrastructure: smaller bodies, endothermy, enhanced metabolic
   efficiency. The exogenous shock cleared the niche above them.
   They crossed.
\item \textbf{Crocodilians.} Stayed on Strategy~1 in a niche where
   the crossover never arrives. Cold-blooded, low metabolic rate,
   indeterminate growth. The domain stays small enough that $O(n^2)$
   works. They have not changed in 200 million years because the
   niche has not changed.
\item \textbf{Non-avian dinosaurs.} Too large for Strategy~1, had not
   built Strategy~2. Sitting in the middle of the sigmoid when the
   shock arrived. Extinct.
\end{enumerate}

Current language models are building feathers. The question is not
whether the transition happens---the scaling math guarantees it---but
which systems are birds and which are crocodilians. Both survive.
LLMs handling bounded-scope tasks (conversation, creative writing,
single-document analysis) may never hit the crossover. They are the
crocodile niche: the domain is small enough. LLMs attempting
open-ended reasoning over large knowledge bases will hit the
crossover. They must build the graph or hit the wall.

The systems that will not survive are the ones scaling the
dinosaur: making the context window bigger, adding more parameters,
paying the quadratic cost---too large for Strategy~1, not yet
building Strategy~2. The middle of the sigmoid is where extinction
happens.

% ------------------------------------------------------------------
\section{Convergence}
\label{sec:convergence}
% ------------------------------------------------------------------

The genome and the software industry arrived at the same architecture:

\begin{center}
\small
\begin{tabular}{lll}
\toprule
\textbf{Architectural pattern} & \textbf{Genome} & \textbf{Software} \\
\midrule
Directed graph with typed edges & Regulatory network & Knowledge graph \\
Hierarchical error correction & 7 repair layers & Test/staging/prod \\
Write-ahead log & Epigenetic marks & Transaction log \\
Selective decompression & Cell-type-specific expression & Lazy loading \\
Encoding permanence & DNA $>$ epigenetics $>$ protein & Code $>$ config $>$ runtime \\
Autoscaling with policy & Mitosis + contact inhibition & Container orchestration \\
Horizontal + vertical scaling & Cell division + specialization & Replicas + microservices \\
Orchestrator as SPOF & TP53 & Load balancer \\
\bottomrule
\end{tabular}
\end{center}

This convergence is not metaphor. It is not analogy. It is the
consequence of a mathematical constraint. Papers~1--4 prove that
any persistent system under bounded context must maintain these
structures. The genome and the software industry are both persistent
systems under bounded context. They found the same solution because
it is the only solution.

The timescales differ by eight orders of magnitude (3.8 billion years
versus 60 years). The substrates differ entirely (nucleotides versus
silicon). The selection pressures differ (survival versus profit).
But the constraint is the same. And the constraint determines the
architecture.

% ------------------------------------------------------------------
\section{What the Series Claims}
\label{sec:claims}
% ------------------------------------------------------------------

Taken together, the eleven papers make the following claims:

\begin{enumerate}
\item Any persistent system under bounded active context must
   maintain a directed graph with typed edges (Paper~1).

\item Such systems converge to escalation architectures under
   survival pressure (Paper~2).

\item Knowledge and action are inseparable under this constraint
   (Paper~3).

\item Information encodes at permanence levels proportional to
   confidence (Paper~4).

\item Cancer follows the topology of the biological knowledge graph:
   channels, positions, and co-occurrence patterns carry independent
   survival signal (Papers~5--8).

\item The genome is an instance of this architecture: DNA as commit
   log, epigenetics as write-ahead log, repair layers as governance
   scaling logarithmically with cell count (Paper~10).

\item The forward ratchet (evolution) and the backward ratchet
   (cancer) are opposing failures of the same governance mechanism.
   The sigmoid crossover between scaling classes determines when each
   occurs (this paper).

\item Software systems under bounded context discover the same
   architecture independently (Paper~9).

\item The convergence is not analogy. It is mathematical necessity.
\end{enumerate}

\subsection{What this series does not claim}

It does not claim that biology is a computer. It claims that
biology and computers are both subject to the same information-theoretic
constraint, and that the constraint admits one family of solutions.

It does not claim that cancer can be ``solved'' by better engineering.
It claims that the structural nature of the problem is identical to
problems that distributed systems engineering has formalized, and that
the formalism may yield insight.

It does not claim that evolution is ``designed.'' It claims that
selection under bounded context is a sufficient condition for the
emergence of graph structure, and that no designer is needed for
the architecture to converge.

% ------------------------------------------------------------------
\section{The Dangerous Property of the Inferior Scaling Class}
\label{sec:dangerous}
% ------------------------------------------------------------------

The dangerous property of $O(n^2)$ is not that it fails. It is that
it succeeds.

At small scale, quadratic search is faster than logarithmic retrieval.
The constant overhead of maintaining an index exceeds the cost of
brute-force search. A flat transformer outperforms a structured
retriever when the corpus is small. A bacterial colony outperforms a
multicellular organism when the environment is simple. A startup with
no process outperforms a company with governance when the team is
five people.

The success is real. The success is measurable. The success is
temporary.

The crossover is invisible from below. Not because it is hidden.
Because the sigmoid is so large that the curve looks flat from where
you stand. The system that is winning at $n = 100$ sees a log curve.
It extrapolates and sees more log curve. The inflection at
$n = 10{,}000$ is real, but it is beyond the horizon. The scaling
exponent is a fact about the function, not about the current input.
By the time the curvature becomes visible, the system is past it.

Cancer follows this curve exactly. A cell with one governance layer
lost is indistinguishable from its neighbors. Two layers lost, still
functional. Three layers lost, measurably different but not yet
dangerous. Each step looks like the last. The curve looks flat. Then
four layers lost---past the inflection---and the cell is replicating
without constraint. Cancer appears sudden because the sigmoid was
too big to see from inside it.

Dinosaurs followed this curve. 165 million years of dominance.
Strategy~1 at enormous scale: large bodies, high metabolic cost,
brute-force resource acquisition. The sigmoid was always there.
They could not see it because the flat part lasted longer than
most lineages exist. When the exogenous shock arrived (the K-Pg
asteroid), three outcomes followed:

\begin{itemize}
\item The lineages that had already begun building $O(\log n)$
   infrastructure---smaller bodies, endothermy, metabolic
   efficiency---survived and radiated. These became birds.
\item The lineages that remained on Strategy~1 in a niche where
   the crossover never arrives---cold-blooded, low metabolic rate,
   small effective $n$---survived unchanged. These are crocodilians.
\item The lineages in between---too large for Strategy~1, too
   invested in the old architecture to build Strategy~2---went
   extinct. The middle of the sigmoid is where extinction happens.
\end{itemize}

This is the prediction the series makes about artificial intelligence.
The current generation of language models operates on the flat part of
the curve. Context windows are growing. Performance is improving.
The success is real. The curve looks like a log from here. It is not.

Some systems will cross: build graph structure, governance layers,
retrieval infrastructure. Birds. Some will find the niche where
quadratic attention is sufficient---bounded tasks, single-document
scope, conversational interaction---and remain there indefinitely.
Crocodilians. Both survive.

The systems that will not survive are the ones scaling the dinosaur:
making the context window bigger, adding parameters, paying the
quadratic cost at increasing scale. Too large for Strategy~1. Not
building Strategy~2. The middle of the sigmoid. The scaling math
does not negotiate.

% ------------------------------------------------------------------
\section{Conclusion}
\label{sec:conclusion}
% ------------------------------------------------------------------

The ratchet clicks forward when a system acquires the governance to
handle more complexity. It clicks backward when governance degrades
and the system can no longer maintain its own structure. Forward is
evolution. Backward is cancer. The mechanism is governance under
bounded context. The shape is a sigmoid. The constraint is
mathematical.

The genome and the software industry are not analogous systems. They
are the same system, implemented on different substrates, operating
at different timescales, subject to the same constraint. The
architecture they share was not borrowed from one domain to the other.
It was derived independently, because it is the only architecture
that satisfies the constraint.

Active context is finite. Everything else follows.

\section*{Acknowledgments}

The author used Claude (Anthropic) for drafting assistance and
\LaTeX{} preparation. All theoretical content and arguments are the
author's own.

\bibliographystyle{plain}
\begin{thebibliography}{99}

\bibitem{McCarthy2026a}
P.~D. McCarthy.
\newblock Graph structure is necessary for information preservation
under bounded context.
\newblock \emph{SSRN preprint}, 2026.

\bibitem{McCarthy2026j}
P.~D. McCarthy.
\newblock The genome as knowledge graph: governance, scaling, and the
architecture of multicellular life.
\newblock \emph{SSRN preprint}, 2026.

\bibitem{Feinberg2006}
A.~P. Feinberg, R.~Ohlsson, and S.~Henikoff.
\newblock The epigenetic progenitor origin of human cancer.
\newblock \emph{Nature Reviews Genetics}, 7(1):21--33, 2006.

\bibitem{Jaiswal2014}
S.~Jaiswal, P.~Fontanillas, J.~Flannick, et~al.
\newblock Age-related clonal hematopoiesis associated with adverse
outcomes.
\newblock \emph{New England Journal of Medicine}, 371(26):2488--2498, 2014.

\bibitem{Slaughter1953}
D.~P. Slaughter, H.~W. Southwick, and W.~Smejkal.
\newblock ``Field cancerization'' in oral stratified squamous epithelium.
\newblock \emph{Cancer}, 6(5):963--968, 1953.

\bibitem{Bianconi2013}
E.~Bianconi, A.~Piovesan, F.~Facchin, et~al.
\newblock An estimation of the number of cells in the human body.
\newblock \emph{Annals of Human Biology}, 40(6):463--471, 2013.

\bibitem{Heard2014}
E.~Heard and R.~A. Martienssen.
\newblock Transgenerational epigenetic inheritance: myths and mechanisms.
\newblock \emph{Cell}, 157(1):95--109, 2014.

\end{thebibliography}

\end{document}
